\section{The model}\label{sec:model}

The economy is a small open economy that imports its fossil energy, populated
by heterogeneous households, competitive firms producing a domestic good from
labour, a union that sets wages, a government, and a rest of the world that sets
world prices. With one exception, the environment is the small open economy HANK of
\citet{auclert2023energy}.

The base is the \citet{auclert2023energy} model. We change one thing, in the household
block: the price of energy a household faces is no longer given but depends on a
durable it chooses. Each household holds a brown or a green durable (a gas boiler
or a heat pump, a combustion or an electric vehicle) and pays a one-time cost
$\psi_g$ to switch from brown to green, after which it buys energy at the lower
green price. The change is confined to the household's problem, but it reaches the
rest of the model through the price it controls: the energy price sits inside the
consumption bundle and the budget constraint. Adoption is an intertemporal choice:
the household weighs the upfront cost $\psi_g$ against the stream of energy savings
the green durable delivers over its life, and pays $\psi_g$ out of current resources
against the borrowing constraint; only households with enough liquidity can afford
it. A shock or a policy that moves the brown energy price moves this decision, and
with it the whole household block. Everything outside that, the CES nest, sticky
import prices, the wage Phillips curve, the mutual fund, and the
constant-real-rate rule, follows AMRS.

Section~\ref{sec:hh} presents the household block.
Sections~\ref{sec:firms}--\ref{sec:gov} lay out the production side, the price
system, the open-economy accounts, and the government, which are standard.
Section~\ref{sec:eqm} defines equilibrium, and Section~\ref{sec:ss} the steady
state; the two crisis experiments are defined in Section~\ref{sec:shocks}. Time is discrete and quarterly. A bar denotes
a steady-state value; a date subscript indexes the period, with $x_{t-1}$ a lag and
$x_{t+1}$ an expected lead. For the durable, $d_{t-1}$ is the technology chosen last period, $d^-_t$ the
technology held at the start of $t$ after breakdown, and $d_t$ the technology chosen. All retail
prices are expressed relative to the CPI $P$ and carry the subscript ${}_{/P}$.

\subsection{Households}\label{sec:hh}

There are three permanent discount-factor types $i\in\{0,1,2\}$ of equal mass, with
discount factors $\beta_i$ chosen in Section~\ref{sec:cal} to match aggregate MPCs.
A household is described by three state variables: its assets $a$, its idiosyncratic
labour productivity $e$, and the durable technology $d\in\{B, G\}$ it operates.
Productivity follows a discretised log-AR(1), $\log e_{t+1}=\rho_e\log e_t+\epsilon_{t+1}$, with a
mean-one grid and transition matrix $\Pi$, as in \citet{auclert2023energy}. The technology is
the new object, and because the paper concerns a decision over it, we fix when,
within a period, the household holds it, loses it, and chooses it.

\subsubsection{The technology and the timing of the choice}\label{sec:tech}

A household enters the period operating the technology it chose last period, which we
denote $d_{t-1}$. Three events then occur, in this order, before it consumes.

\begin{enumerate}
\item \emph{Breakdown.} A green durable fails with probability $\delta_g$ and reverts
to brown; a brown durable is unaffected. We write $d^-_t$ for the technology the
household holds \emph{after} this draw: the technology it brings to the decision.
Hence $d^-_t=B$ whenever $d_{t-1}=B$, while $d^-_t=G$ with probability $1-\delta_g$ when
$d_{t-1}=G$.
\item \emph{Productivity.} The household draws its labour productivity $e_t$.
\item \emph{The decision.} The household chooses, in a single problem, the
technology $d_t$ it will operate together with its consumption and saving. A household
holding brown ($d^-_t=B$) may keep brown or pay the switching cost $\psi_g$ to adopt
green; a household holding green ($d^-_t=G$) operates it, since running an owned green
durable costs nothing and reverting to brown is never chosen.
\end{enumerate}

No information arrives between the productivity draw and the decision, so the
household settles technology, consumption, and saving \emph{jointly}
(Section~\ref{sec:problem}). The technology it operates is what it carries
forward, so $d_t$ becomes the technology held on entry next period; the within-period
sequence is
%
\[
\underbrace{d_{t-1}}_{\text{held on entry}}
\;\xrightarrow{\ \text{breakdown}\ }\;
\underbrace{d^-_t}_{\text{held at the decision}}
\;\xrightarrow{\ \text{choice}\ }\;
\underbrace{d_t}_{\text{operated, carried forward}}.
\]
%
Two binary objects thus describe the technology: the holding at the decision, $d^-_t$,
and the operated technology, $d_t$. The switching cost is paid exactly when a brown
holder adopts, that is when $d^-_t=B$ and $d_t=G$; we write $\I[d^-_t=B, d_t=G]$ for a switcher
and $\I[d_t=G]$ for a household operating green. The feasible choices are
%
\[
\mathcal F(d^-_t)=
\begin{cases}
\{B,\,G\} & d^-_t=B \quad(\text{keep brown, or pay $\psi_g$ to adopt}),\\[2pt]
\{G\} & d^-_t=G \quad(\text{operate the green durable}).
\end{cases}
\]

Two features follow directly. First, the technology is
persistent but asymmetric: brown is absorbing, and green is lost only through
breakdown. Second, the adoption decision is an option held by the brown
population alone: a green household has nothing to choose.

A constant green share therefore requires a standing flow of switchers to replace the
durables that fail. In steady state this flow is
%
\begin{equation}
\bar D^{\mathrm{sw}}=\delta_g\,\bar\DG,
\label{eq:switch-flow}
\end{equation}
%
with $\DG$ and $\dsw$ the population shares of green households and of switchers. At
the calibrated $\bar\DG=5\%$ and a five-year durable life, this flow is about a
quarter of one percent of households per quarter: adoption is a rare, lumpy event,
so the object to identify is its elasticity rather than its level
(Section~\ref{sec:taste_id}, Appendix~\ref{app:taste}).

\subsubsection{Preferences and the cost of living}\label{sec:prefs}

Labour supply is determined at the union level (Section~\ref{sec:firms}), so flow
utility is CRRA over a consumption bundle $c$ alone,
%
\begin{equation}
u(c)=\frac{c^{\,1-\sigma}}{1-\sigma},
\qquad
u(c)=\log c \ \text{ for } \sigma=1,
\label{eq:felicity}
\end{equation}
%
with $\sigma$ the inverse elasticity of intertemporal substitution, as in \citet{auclert2023energy}. The bundle is the
\citet{auclert2023energy} CES nest of energy against a home/foreign goods composite, with
energy weight $\alpha_E$ and a low substitution elasticity $\eta_E$.

The departure from AMRS is that energy comes at a technology-specific price: a
household operating brown buys energy at $\PEB$, one operating green at $\PEG<\PEB$.
The operated technology therefore fixes both the energy price and the exact
cost-of-living index,
%
\begin{equation}
p^E_{d,t}=\PEB[,t]+\I[d=G]\,\bigl(\PEG[,t]-\PEB[,t]\bigr),
\label{eq:pEd}
\end{equation}
\begin{equation}
\frac{p_{G,t}}{p_{B,t}}=\left[\frac{(1-\alpha_E)\,p_{HF/P,t}^{\,1-\eta_E}+\alpha_E\,(\PEG[,t])^{1-\eta_E}}
{(1-\alpha_E)\,p_{HF/P,t}^{\,1-\eta_E}+\alpha_E\,(\PEB[,t])^{1-\eta_E}}\right]^{\frac{1}{1-\eta_E}},
\qquad
\frac{p_{G,t}}{p_{B,t}}=\bigl(\PEG[,t]/\PEB[,t]\bigr)^{\alpha_E}\ \text{ for }\eta_E=1.
\label{eq:pd}
\end{equation}
%
The adoption choice responds to the relative saving $p_G/p_B<1$, the cheaper cost
of living the switching cost buys; at the calibrated steady state $p_G/p_B\approx0.981$.

\subsubsection{The household's problem}\label{sec:problem}

A household of type $i$ reaches period $t$ with entry assets $a_{t-1}$ and
productivity $e_t$, holding the post-breakdown durable $d^-_t$, and chooses its
technology $d_t$, consumption $c_t$, and saving $a_t$. Its resources are three.
Assets earn the ex-post real return $r_t$,
the return on the mutual fund of Section~\ref{sec:firms}. This return equals the
lagged ex-ante rate $r^{\mathrm{ante}}_{t-1}$ at every date but the first; at
date~0 the revaluation of the fund's equity drives the two apart
(equation~\eqref{eq:fund}), and the household bears that surprise. Labour income is $e_t\,Z_t$, where $Z_t\equiv(1-\tau^L_t)w_t n_t$
is aggregate after-tax labour income distributed in proportion to productivity, as in
\citet{auclert2023energy}; the household does not choose hours, which the union sets
(Section~\ref{sec:firms}). The government may pay a transfer $T_t$, specified below. On
the spending side, the household buys the consumption bundle at its
technology-specific cost of living $p_{d,t}$, saves $a_t\ge0$ subject to the borrowing
constraint, and, if it is a brown holder adopting green this period, pays the
switching cost $\psi_g$, once. The budget constraint collects these
items,
%
\begin{equation}
p_{d,t}\,c_t+a_t+\psi_g\,\I[d^-_t=B,\,d_t=G]=(1+r_t)\,a_{t-1}+e_t\,Z_t+T_t.
\label{eq:coh}
\end{equation}

The switching cost is paid exactly once, in the period a brown holder adopts, and
it buys the permanently cheaper bundle $p_G<1$ for as long as the durable
survives: an upfront, lumpy outlay against a gradual operating saving. Because the outlay is financed out of current
resources against the borrowing constraint, wealth governs who can undertake it,
which is why adoption concentrates among the unconstrained
(Section~\ref{sec:ss_adoption}).

The transfer carries ex-post fiscal policy,
%
\begin{equation}
T_t=T^{\mathrm{unt}}_t+ins_E\,(\pE_t-\bar p^E)\,\bar c^{\,B}_{E,i}(e, a),
\label{eq:transfer-hh}
\end{equation}
%
an untargeted lump sum $T^{\mathrm{unt}}$ plus a targeted (Slutsky) transfer that
scales, with insurance proportion $ins_E$, the crisis rise in the energy price by the
household's
\emph{pre-crisis} brown-energy demand $\bar c^{\,B}_{E,i}(e, a)$; both are zero in the
baseline. One design feature matters throughout: the transfer conditions on
steady-state energy use, never on the current technology. A household that switches
forfeits nothing, so the transfer leaves the adoption margin undistorted, in
contrast to the price cap of Section~\ref{sec:prices}, which acts on $\PEB$ itself.

The household chooses its technology, consumption, and saving \emph{together}, in a
single problem:
%
\begin{equation}
\max_{d_t\in\mathcal F(d^-_t)}\;\;\max_{c_t,\;a_t\ge0}\;
\Bigl\{\,u(c_t)+\beta_i\,\E_t\bigl[\bar V_{t+1}(d^-_{t+1},e_{t+1},a_t)\bigr]\Bigr\}
\quad\text{subject to \eqref{eq:coh}},
\label{eq:joint}
\end{equation}
%
with $\bar V_{t+1}$ the continuation value (defined in \eqref{eq:logit}) and next
period's holding $d^-_{t+1}$ given by breakdown applied to the operated technology $d_t$. The nested
maxima do not describe a temporal sequence, since nothing is realised between
them; they are the structure of a discrete--continuous choice: for each feasible technology the
household solves the continuous consumption-saving problem, obtaining the
technology-conditional value $V_t(d_t\mid d^-_t,e_t, a_{t-1})$, in which the origin $d^-_t$ enters
only through the switching cost; it then selects the technology
(Section~\ref{sec:adoption}). Only three
continuous problems are ever relevant: a brown household ($p_B$, no switching
cost), a green incumbent ($p_G$, no switching cost), and a green adopter ($p_G$,
resources reduced by $\psi_g$). Each satisfies the usual Euler and envelope
conditions,
%
\begin{equation}
u'(c_t)=\beta_i\,p_{d,t}\,\E_t\bigl[\bar V_{a,t+1}(d^-_{t+1},e_{t+1},a_t)\bigr],
\qquad
V_{a,t}=(1+r_t)\,u'(c_t)/p_{d,t},
\label{eq:euler}
\end{equation}
%
with $u'(c_t)=c_t^{-\sigma}$; a green household's cheaper bundle raises the
consumption value of every unit of wealth. CES demands split by the operated
technology,
%
\begin{equation}
c_{E,t}(d)=\alpha_E\bigl(p^E_{d,t}/p_{d,t}\bigr)^{-\eta_E}c_t,
\qquad
c_{HF,t}(d)=(1-\alpha_E)\bigl(p_{HF/P,t}/p_{d,t}\bigr)^{-\eta_E}c_t,
\label{eq:hh-demands}
\end{equation}
%
and brown versus green energy use follow the technology,
$c^B_{E,t}=c_{E,t}(d_t)\,(1-\I[d_t=G])$ and $c^G_{E,t}=c_{E,t}(d_t)\,\I[d_t=G]$.

\subsubsection{The technology choice}\label{sec:adoption}

The outer maximum in \eqref{eq:joint} selects the technology, completing the joint
problem. A household holding $d^-_t$ compares the technology-conditional values
$V_t(d_t\mid d^-_t,\cdot)$; with an i.i.d.\ extreme-value taste shock of scale
$\sigma_\varepsilon$ on each option, the choice is logit and the household's
continuation value $\bar V_t$ is the log-sum-exp,
%
\begin{equation}
\begin{aligned}
\Pr\nolimits_t(d_t\mid d^-_t,e_t, a_{t-1})&=
\frac{\exp\{V_t(d_t\mid d^-_t,e_t, a_{t-1})/\sigma_\varepsilon\}}
{\sum_{d'\in\mathcal F(d^-_t)}\exp\{V_t(d'\mid d^-_t,e_t, a_{t-1})/\sigma_\varepsilon\}},\\[4pt]
\bar V_t(d^-_t,e_t, a_{t-1})&=\sigma_\varepsilon\log\!\!\sum_{d'\in\mathcal F(d^-_t)}\!\!
\exp\bigl(V_t(d'\mid d^-_t,e_t, a_{t-1})/\sigma_\varepsilon\bigr).
\end{aligned}
\label{eq:logit}
\end{equation}
%
For a green holder the feasible set is a singleton, so $\bar V_t(G,\cdot)=V_t(G\mid G,\cdot)$
and there is nothing to decide. The only genuine choice is the brown holder's: a
two-way comparison of staying brown against paying $\psi_g$ to adopt green. That cost
already sits inside the adopter's budget constraint \eqref{eq:coh}, so the option values
price it automatically.

The economics of the margin is simplest in steady state. The marginal switcher
equates the premium to the discounted operating-cost saving over the durable's
expected life,
%
\begin{equation}
\psi_g\;\approx\;\frac{\bigl(\bar P^E_B-\bar P^E_G\bigr)\,\bar C^G_E}{r+\delta_g},
\label{eq:margin}
\end{equation}
%
which ties the adoption threshold to the delivered cost ratio
$\varrho=\bar P^E_G/\bar P^E_B$. Every experiment in the paper acts on this one
comparison. An energy crisis raises $\PEB$ and lifts the return to switching; a
retail price cap holds $\PEB$ at its pre-crisis level and removes that return
entirely; a transfer supports income while leaving $\PEB$, and hence
\eqref{eq:margin}, free; a permanent carbon tax raises the brown price in steady
state, greening the economy before any crisis arrives.

Away from steady state, the responsiveness of the margin is governed by
$\sigma_\varepsilon$. The taste shocks are, in the tradition of \citet{iskhakov2017endogenous},
the smoothing that makes the discrete choice differentiable; at a rare margin they are
also the structural semi-elasticity of switching, so $\sigma_\varepsilon$ cannot be
treated as a nuisance parameter and set small. Section~\ref{sec:taste_id} identifies
it from quasi-experimental adoption elasticities, and Appendix~\ref{app:taste}
develops the smoothing--elasticity duality.

\subsubsection{Aggregation}\label{sec:agg2}

For each type $i$, aggregates integrate the household policies over the stationary
distribution $\Lambda_i$, and the economy averages the three types equally,
%
\begin{equation}
X_t=\tfrac13\textstyle\sum_i X_{i,t},
\qquad
X\in\{C,\;A,\;\DG,\;\dsw,\;C^B_E,\;C^G_E,\;C_{HF},\dots\}.
\label{eq:agg}
\end{equation}
%
Aggregate energy demand is $c_{E,t}=C^B_{E,t}+C^G_{E,t}$, and home- and foreign-good demands apply
the common inner nest to the composite,
$c_{H,t}=(1-\alpha_F)\,p_{H/HF,t}^{-\eta}\,C_{HF,t}$ and
$c_{F,t}=\alpha_F\,p_{F/HF,t}^{-\eta}\,C_{HF,t}$. The two aggregates the policy experiments
track are the green share $\DG$, the population share operating green, and the
switching flow $\dsw$, the share adopting green each period.

\subsection{Production and wage setting}\label{sec:firms}

\paragraph{Production and profits.}
Competitive firms produce the domestic final good from labour with productivity
$A^N$, sold at a constant gross markup $\mu$, so output is $y=A^N n$ and after-tax
aggregate labour income is $Z=(1-\tau^L)wn$, with $n$ hours and $w$ the real wage. Monopoly profits are capitalised into a mutual fund that households
hold,
%
\begin{equation}
D_t=(\mu-1)\,w_t n_t,
\qquad
J_t=D_t+\frac{J_{t+1}}{1+r^{\mathrm{ante}}_t},
\qquad
j_t=\frac{J_{t+1}}{1+r^{\mathrm{ante}}_t},
\qquad
1+r_t=\frac{J_t}{j_{t-1}},
\label{eq:fund}
\end{equation}
%
following the notation of \citet{auclert2023energy}: $r^{\mathrm{ante}}_t$ is the ex-ante real
rate at which the risk-neutral fund prices its claims (equation~\eqref{eq:mp}), $j$ is
the ex-dividend share price, and $r$ is the fund's ex-post return: the return
households earn in \eqref{eq:coh}. Because the fund equates ex-ante returns across
assets, the ex-post return equals the lagged ex-ante rate at every date but the
first, $r_t=r^{\mathrm{ante}}_{t-1}$ for $t\ge1$; only the impact return $r_0$ differs, because the date-0 revaluation of the
fund's equity feeds the revaluation term in \eqref{eq:nfa}. Importers are foreign-owned and the energy
endowment is not held domestically, so domestic firms are the only domestically held
risky claim.

\paragraph{Wage setting.}
A union sets the nominal wage under Calvo frictions $\theta_w$, trading off
households' marginal rate of substitution against the after-tax wage, with a
real-wage-stabilisation term (weight $\zeta_{BG}$) as in
\citet{auclert2023energy}. With slope $\kappa_w=(1-\theta_w)(1-\beta\theta_w)/\theta_w$,
inverse Frisch elasticity $\varphi$, and labour-disutility constant $v_\varphi$,
wage inflation follows
%
\begin{equation}
\pi^w_t(1+\pi^w_t)
= \beta\,\pi^w_{t+1}(1+\pi^w_{t+1})
+ \frac{\kappa_w}{v_\varphi\,n_t^{\varphi}}
\biggl[v_\varphi\,n_t^{\varphi}
- \frac{(1-\tau^L_t)w_t\,C_t^{-\sigma}}{\mu}
- \zeta_{BG}\frac{(w_t-\bar w)\,w_t\,\bar C^{-\sigma}\bar n}{\mu\,n_t\,\bar w}\biggr],
\label{eq:wpc}
\end{equation}
%
and the real wage evolves as $w_t=(1+\pi^w_t)\,w_{t-1}/(1+\pi_t)$.

\subsection{The price system}\label{sec:prices}

The paper's mechanism runs entirely through relative prices, so we lay the price
system out in one place. Table~\ref{tab:prices} collects it: world prices at the
top, a wholesale layer carrying the nominal rigidities, a retail layer where
policy acts, and the CPI that aggregates them.

\begin{table}[H]
\centering
\figtitle{The price system}
\small
\begin{tabular}{lll}
\toprule
Symbol & Object & Determined by \\
\midrule
$P^{E*},\,P^{*}_F,\,C^{*},\,r^{*}$ & world energy and goods prices, demand, rate &
exogenous\textsuperscript{a} \\
$\pE$ & wholesale brown energy price & NKPC \eqref{eq:pE-nkpc} \\
$\PEB$ & brown households' energy price (tax-inclusive) & cap rule and $\tau_b$, eq.~\eqref{eq:cap} \\
$\PEG$ & green households' energy price (tax-inclusive) & world price and $\tau_g$, eq.~\eqref{eq:pEG2} \\
$\PEBpre,\,\PEGpre$ & pre-tax energy prices (BoP, carbon account) & eqs.~\eqref{eq:cap}--\eqref{eq:pEG2} \\
$p_{H/P}$ & domestic good & markup pricing \eqref{eq:pH} \\
$p_{F/P}$ & imported consumption good & NKPC \eqref{eq:pF} \\
$p_{HF/P},\,p_{H/HF}$ & goods composites & outer nest \eqref{eq:outernest} \\
$p_B,\,p_G$ & brown/green cost of living & bundle nests \eqref{eq:innernest};
saving $p_G/p_B$ \eqref{eq:pd} \\
$Q$ & real exchange rate & UIP \eqref{eq:uip} \\
\bottomrule
\end{tabular}
\label{tab:prices}
\fignotes{All retail prices are relative to the CPI. \textsuperscript{a}Under the
fixed-quantity closure $P^{E*}$ is the price that clears the fixed regional quantity
(Section~\ref{sec:eqm}).}
\end{table}

\paragraph{The domestic good.}
Marginal-cost pricing at the constant markup gives
%
\begin{equation}
p_{H/P,t}=\frac{\mu}{A^N}\,w_t.
\label{eq:pH}
\end{equation}

\paragraph{Imported goods and wholesale energy.}
Retailers of the imported good and of energy reset prices under Calvo frictions
$\theta_F$ and $\theta_E$, so world price movements pass through to domestic
retail prices only gradually. With the usual slopes $\kappa_F$ and $\kappa_E$,
%
\begin{equation}
\kappa_F\Bigl(\frac{Q_t P^{*}_F}{p_{F/P,t}}-1\Bigr)
+ \beta_F\,\pi^F_{t+1}(1+\pi^F_{t+1}) - \pi^F_t(1+\pi^F_t)=0,
\label{eq:pF}
\end{equation}
\begin{equation}
\kappa_E\Bigl(\frac{Q_t P^{E*}_t}{\pE_t}-1\Bigr)
+ \beta_E\,\pi^E_{t+1}(1+\pi^E_{t+1}) - \pi^E_t(1+\pi^E_t)=0.
\label{eq:pE-nkpc}
\end{equation}
%
The margins these retailers earn out of steady state accrue abroad (the importers
are foreign-owned, as in \citealp{auclert2023energy}); they matter below only because the
import bill is valued at retail rather than world prices.

\paragraph{Retail energy prices.}
Each technology's energy price is built in two steps: a pre-tax price, which
values the resource flows in the balance of payments and the carbon account, and
the tax-inclusive price the household pays. For brown energy the pre-tax price
applies the cap to the wholesale price, and the carbon tax scales it,
%
\begin{equation}
\PEBpre[,t]=(1-\tau^E)\,\pE_t+\tau^E\,\bar p^E,
\qquad
\PEB[,t]=\PEBpre[,t]\,(1+\tau_b),
\label{eq:cap}
\end{equation}
%
so $\tau^E$ governs pass-through continuously: at $\tau^E=0$ households face the
wholesale price, at $\tau^E=1$ the pre-tax brown price is frozen at its
pre-crisis level: the full shield. Green energy is imported at the constant
world price $P^{E*}_G=\varrho\,\bar p^E$, $\varrho<1$, and carries its own
carbon tax $\tau_g$,
%
\begin{equation}
\PEGpre[,t]=Q_t\,P^{E*}_G,
\qquad
\PEG[,t]=\PEGpre[,t]\,(1+\tau_g),
\label{eq:pEG2}
\end{equation}
%
so adopters are fully insulated from the fossil shock: the real exchange rate
does not move under either shock, and $\PEG$ stays at $\varrho\,\bar p^E$.

The household side of the model reads only the tax-inclusive prices: the bundle
nests \eqref{eq:innernest}, the budget \eqref{eq:coh} and hence the adoption
margin \eqref{eq:margin} all price energy at $\PEB$ and $\PEG$. The pre-tax
prices reappear where resources cross the border, in the balance of payments
(Section~\ref{sec:bop}), and in the carbon account of Section~\ref{sec:gov}. In
the baseline $\tau_b=\tau_g=0$ and the two coincide; under the carbon-tax
experiments of Section~\ref{sec:exante} they diverge, and the distinction is what
gives the tax real incidence: it raises the price the adoption margin compares
without changing the resource cost of energy to the country.

\paragraph{The consumer price index.}
Each technology has its own consumption bundle. Writing $p_B,p_G$ for the brown
and green bundle prices relative to the CPI, the inner (energy versus goods) CES
nests are
%
\begin{equation}
p_{d,t}^{\,1-\eta_E}=(1-\alpha_E)\,p_{HF/P,t}^{\,1-\eta_E}+\alpha_E\,(P^E_{d,t})^{1-\eta_E},
\qquad d\in\{B,G\},
\label{eq:innernest}
\end{equation}
%
with the outer (home versus foreign) nest
%
\begin{equation}
(1-\alpha_F)\,p_{H/HF,t}^{\,1-\eta}+\alpha_F\,p_{F/HF,t}^{\,1-\eta}=1.
\label{eq:outernest}
\end{equation}
%
The CPI is the expenditure-weighted index of the two bundles: $P$ is defined by the
requirement that nominal consumption expenditure deflated by $P$ equal real
consumption, $\int p_{d(i)}\,c(i)\,\mathrm{d}i=\int c(i)\,\mathrm{d}i$, so that
$\sum_d \omega_{d,t}\,p_{d,t}=1$ with $\omega_{d,t}$ the consumption share of state $d$. As
green diffuses the energy leg is reweighted from brown toward green through
$\omega_{G,t}=\DG_t$, so $P$ is the true cost-of-living index the monetary rule reads, with no measurement
gap between the price index and the basket households actually buy.

\subsection{The open economy}\label{sec:open}\label{sec:bop}

\paragraph{Foreign demand and the exchange rate.}
Foreign demand for the home good is $c^{*}_H=\alpha^{*}(p^{*}_H)^{-\gamma}C^{*}$,
with trade elasticity $\gamma$, and the real exchange rate is pinned by uncovered
interest parity,
%
\begin{equation}
\frac{Q_t}{Q_{t+1}}\bigl(1+r^{\mathrm{ante}}_t\bigr)=1+r^{*}.
\label{eq:uip}
\end{equation}

\paragraph{World energy supply.}
World supply is a fixed endowment $E^{*}$ plus intertemporal stock
arbitrage with cost $\Gamma$,
%
\begin{equation}
E^{\mathrm{stock}}_t=\frac{P^{E*}_{t+1}/(1+r^{*})-P^{E*}_t}{\Gamma},
\qquad
E^{s}_t=E^{*}_t
+\bigl(E^{\mathrm{stock}}_{t-1}-E^{\mathrm{stock}}_t\bigr).
\label{eq:energy-supply}
\end{equation}
%
The home economy owns none of the endowment, being a pure importer, so energy
rents never enter domestic wealth.

\paragraph{The balance of payments and the adoption flows.}
The adoption margin adds exactly two terms to the current account, and they pull
in opposite directions. The switching cost is a resource outflow, booked as an
import $M^{\mathrm{sw}}$: it is front-loaded, paid once per switcher, and scales with
the switching volume. The adopters' saved fossil energy is a resource inflow, booked
as an import saving $S^{E}$: it accrues gradually, over the life of each green
durable,
%
\begin{equation}
M^{\mathrm{sw}}_t=\psi_g\,\dsw_t,
\qquad
S^{E}_t=\bigl(\PEBpre-\PEGpre\bigr)\,C^{G}_{E,t}.
\label{eq:adoption-flows}
\end{equation}
%
The energy saving is valued at pre-tax prices: the carbon taxes are transfers
between domestic households and the government, not resource flows to the rest
of the world. Their revenue,
$R^{\mathrm{carbon}}=\tau_b\,\PEBpre\,C^B_E+\tau_g\,\PEGpre\,C^G_E$, accrues to
the government budget (Section~\ref{sec:gov}) and never enters net exports.
Net exports collect trade in goods and energy together with the two adoption
flows,
%
\begin{equation}
nx_t = Q_t\,p^{*}_H c^{*}_{H,t} - p_{F/P,t}\,c_{F,t} - \pE_t\,c_{E,t}
- M^{\mathrm{sw}}_t + S^{E}_t,
\label{eq:nx}
\end{equation}
%
and net foreign assets accumulate at the ex-ante rate, with one additional term,
%
\begin{equation}
nfa_t = nx_t + (r_t-r^{\mathrm{ante}}_{t-1})\bigl(A_{t-1}-j_{t-1}\bigr)
+ (1+r^{\mathrm{ante}}_{t-1})\,nfa_{t-1}.
\label{eq:nfa}
\end{equation}
%
The middle term is a revaluation. Households earn the
ex-post return $r$ on their wealth, while the country's accounts accrue at the
promised ex-ante rate $r^{\mathrm{ante}}_{t-1}$; the two differ only at date~0, when
the shock revalues the equity of the mutual fund (equation~\eqref{eq:fund}).
Applying the surprise $r_t-r^{\mathrm{ante}}_{t-1}$ to households' wealth net of the
domestically held fund, $A_{t-1}-j_{t-1}$, records the one-off capital
gain or loss on the country's external claims. At date~0, the only date the term is
active, $A_{t-1}-j_{t-1}$ equals the net foreign position itself, since
government debt is zero at the pre-crisis steady state; from date~1 onward
$r_t=r^{\mathrm{ante}}_{t-1}$ and the term vanishes.

Booking the adoption outlay as an import is a convention rather than a result. The
baseline sets the import share of the switching bundle to one,
$\alpha_F^{switch}=1$, so the full cost $\psi_g\,\dsw$ enters \eqref{eq:nx} as
$M^{\mathrm{sw}}$; Appendix~\ref{app:booking} traces the continuous import share
from one down to zero, where the outlay is spent entirely on the home good. The
convention fixes the sign of the adoption channel's \emph{output} contribution
but none of the policy rankings, which run through relative prices alone.

\subsection{The government}\label{sec:gov}

\paragraph{Monetary policy.}
The nominal rate follows a rule in current and expected inflation, with the
real rate defined by the Fisher equation,
%
\begin{equation}
1+i_t=(1+\bar r^{*})(1+\phi_\pi\pi_t)(1+\phi_{\pi e}\E_t\pi_{t+1})+\varepsilon^m_t,
\qquad
r^{\mathrm{ante}}_t=\frac{1+i_t}{1+\E_t\pi_{t+1}}-1,
\label{eq:mp}
\end{equation}
%
where $\varepsilon^m$ is a monetary shock, zero outside the monetary
experiment of Appendix~\ref{sec:monetary}. The baseline is the constant-real-rate
rule of \citet{auclert2023energy}, $(\phi_\pi,\phi_{\pi e})=(0,1)$. Through \eqref{eq:uip} it
pins the path of the real exchange rate, so a domestic transfer has no
terms-of-trade effect, a property inherited from AMRS rather than assumed anew. A standard
Taylor rule, $(\phi_\pi,\phi_{\pi e})=(1.5,0)$, is reported as a robustness
contrast in Appendix~\ref{sec:monetary}.

\paragraph{Fiscal policy.}
Fiscal policy is the object of Sections~\ref{sec:shield}--\ref{sec:exante}. Four
instruments are deployed once the shock hits and are zero in the baseline; a
fifth, the carbon tax, is permanent and defines the ``prepared'' economy of
Section~\ref{sec:exante}. Table~\ref{tab:instruments} states,
for each, what it acts on, the base its cost scales with, and its budgetary cost.

\begin{table}[H]
\centering
\figtitle{Fiscal instruments}
\small
\begin{tabular}{llll}
\toprule
Instrument & Acts on & Cost scales with & Fiscal cost \\
\midrule
Price cap ($\tau^E$) & brown retail price $\PEB$ & actual energy use & $\tau^E(\pE-\bar p^E)\,c_E$ \\
Slutsky transfer ($ins_E$) & household income & pre-crisis brown use & $ins_E\,(\pE-\bar p^E)\,\bar C^B_E$ \\
Flat transfer ($T^{\mathrm{unt}}$) & household income & none (lump sum) & $T^{\mathrm{unt}}$ \\
Green subsidy ($s_g$) & switching cost $\psi_g$ & number of switchers & $s_g\,\psi_g\,\dsw$ \\
\midrule
Carbon tax ($\tau_b$) & brown price $\PEB$, permanent & brown energy use & rebated (App.~\ref{app:booking}) \\
\bottomrule
\end{tabular}
\label{tab:instruments}
\fignotes{The first four are crisis instruments, zero in the baseline
(Section~\ref{sec:shield}); the carbon tax is permanent and defines
the prepared economy (Section~\ref{sec:exante}). Its revenue is rebated lump sum in
the baseline, with green-subsidy recycling in Appendix~\ref{app:booking}.}
\end{table}

The distinction that matters is between acting on the \emph{price} and acting
on \emph{income}. The price cap holds $\PEB$ down through \eqref{eq:cap}; because
$\PEB$ is the price the adoption margin responds to (equation~\eqref{eq:margin}),
capping the price switches the margin off, and the cap's cost is a marginal
subsidy: it scales with actual crisis energy use. The two transfers instead leave
$\PEB$ free and support income, so the margin survives. They differ in their base: the
\emph{Slutsky} transfer (intensity $ins_E$) is indexed to each household's pre-crisis
brown-energy use, so a switcher keeps the same cheque and the margin is
undistorted, but, as Section~\ref{sec:flat} shows, indexing to energy use is
regressive; the \emph{flat} transfer $T^{\mathrm{unt}}$ is the equal-cost,
unconditional alternative. The \emph{green subsidy} pays a
fraction $s_g$ of the switching cost, forcing adoption directly rather than through
the price. The \emph{carbon tax} raises $\PEB$ permanently, greening the economy
before any crisis; it is the ex-ante counterpart to the crisis instruments and the
subject of Section~\ref{sec:exante}.

All crisis outlays are deficit-financed. Collecting them term by term, with each the
cost line of Table~\ref{tab:instruments},
%
\begin{equation}
X_t=\underbrace{T^{\mathrm{unt}}_t}_{\text{flat transfer}}
+\underbrace{ins_E\,(\pE_t-\bar p^E)\,\bar C^{B}_E}_{\text{Slutsky transfer}}
+\underbrace{\tau^E\,(\pE_t-\bar p^E)\,c_{E,t}}_{\text{price cap}}
+\underbrace{s_g\,\psi_g\,\dsw_t}_{\text{green subsidy}},
\label{eq:outlays}
\end{equation}
%
debt is stabilised by the distortionary labour tax,
%
\begin{equation}
B_t=(1+r^{\mathrm{ante}}_{t-1})B_{t-1}+X_t-\tau^L_t w_t n_t,
\qquad
\tau^L_t=\psi_B\,(B_{t-1}-\bar B),
\label{eq:gbc2}
\end{equation}
%
with no debt in steady state ($\bar B=0$): crisis spending is repaid by raising the
labour tax in proportion to the debt it creates.

\subsection{Equilibrium}\label{sec:eqm}

\paragraph{Numeraire.} The unit of account is the CPI basket, so $p^{\num}=1$ and the
household budget is stated in real terms; $P$ is also the index the monetary rule targets.

\paragraph{Market clearing.}
Goods, assets, and energy clear:
%
\begin{equation}
c_{H,t}+c^{*}_{H,t}-y_t=0,
\qquad
A_t-nfa_t-j_t-B_t=0,
\qquad
\begin{cases}
P^{E*}_t\ \text{exogenous} & \text{(elastic supply)},\\
c_{E,t}=E^{s}_t & \text{(fixed quantity)},
\end{cases}
\label{eq:clearing}
\end{equation}
%
where the asset-market condition is redundant by Walras' law: it holds once goods,
energy, and the government budget balance, so we do not impose it separately. The two
energy closures correspond to the two experiments of Section~\ref{sec:experiments}: a
price-taking economy facing an exogenous world-price path, and a fixed-quantity
economy whose price clears an inelastic supply. Under the fixed-quantity closure the
relevant market is the European one: availability is set region-wide and $P^{E*}$ is
the price that clears it against regional demand.

Let $\Psi_t$ denote the joint distribution over the household state
$(d, e, a)$, durable technology, idiosyncratic productivity, and assets, at date
$t$. Taking as given the sequences of prices, taxes, and transfers, each household
solves the recursive program of Section~\ref{sec:problem}: it chooses nondurable
consumption and savings by \eqref{eq:joint}--\eqref{eq:euler} and, at the durable
stage, switches from brown to green with the logit probability \eqref{eq:logit}.
This program delivers value functions $V_t$ and policy functions
$\{c_t, a_t, P^{sw}_t\}(d, e, a)$. The distribution then evolves as
$\Psi_{t+1}=\Lambda_t(\Psi_t)$, where the law of motion $\Lambda_t$ is induced by
these policies together with the exogenous productivity chain, and aggregates are the
corresponding moments of $\Psi_t$ \eqref{eq:agg}.

\paragraph{Definition (competitive equilibrium).}
Given the world paths $\{P^{E*},r^{*},P^{F*}\}$, the initial distribution $\Psi_0$
and net foreign assets, and the monetary and fiscal rules, a competitive equilibrium
is a sequence of aggregate prices
$\{w,\pi,\pi^{w},\pi^{F},\pi^{E},Q,\pE, r, r^{\mathrm{ante}},p_{H/P}\}_t$, value and
policy functions $\{V_t, c_t, a_t, P^{sw}_t\}$, distributions $\{\Psi_t\}$, aggregate
quantities $\{y, n, c_H, c_F, c_E, C, A\}_t$, technology shares $\{\DG,\dsw\}_t$, and fiscal
variables $\{B,\tau^{L}\}_t$ such that:
\begin{enumerate}[label=(\roman*),leftmargin=2.2em, itemsep=1pt]
\item \emph{(Households)} given prices, taxes, and transfers, the value and policy
functions solve the recursive problem \eqref{eq:joint}--\eqref{eq:euler}, with
brown-to-green switching governed by \eqref{eq:logit}--\eqref{eq:margin};
\item \emph{(Consistency)} $\Psi_t$ evolves from $\Psi_0$ under the law of motion
induced by the policy functions and the productivity process, and the aggregates in
\eqref{eq:agg} are the moments of $\Psi_t$;
\item \emph{(Firms and unions)} the price and wage Phillips curves
\eqref{eq:fund}--\eqref{eq:wpc} and the price system
\eqref{eq:pH}--\eqref{eq:pEG2} hold;
\item \emph{(Government)} the monetary rule \eqref{eq:mp}, the fiscal rule, and the
government budget \eqref{eq:gbc2} hold;
\item \emph{(Open economy)} uncovered interest parity \eqref{eq:uip}, the
balance-of-payments and net-foreign-asset accounting
\eqref{eq:nx}--\eqref{eq:nfa}, and the energy-supply closure
\eqref{eq:energy-supply} hold;
\item \emph{(Clearing)} the goods and energy markets clear \eqref{eq:clearing}; the
asset market clears by Walras' law.
\end{enumerate}

The object of interest is thus a sequence of prices and allocations, together
with a distribution consistent with individual optimisation. The sequence-space Jacobian
method of \citet{Auclert2021a}, applied in Section~\ref{sec:ss}, computes the
linearised transition of this equilibrium.

\subsection{Steady state}\label{sec:ss}

Because the economy is a pure importer, the openness normalisations of \citet{auclert2023energy} reduce to
$\alpha_F=(\alpha-\alpha_E)/(1-\alpha_E)$ and a trade elasticity $\gamma=\eta$ set
to hit an import-demand target $\chi=0.3$, with scale normalisations $y=1$,
$A^N=\mu$, $\alpha^{*}=\alpha$, and $\bar E^{*}=\alpha_E$. The steady
state solves for the labour-disutility shifter $\phi$, the top discount factor
$\beta_{\max}$, output, and the switching cost $\psi_g$, against the wage Phillips
curve, the net-foreign-asset condition, goods clearing, and the green-share target
$\bar\DG=5\%$. The switching cost is thus calibrated to the level of adoption, and
Section~\ref{sec:cal} takes up the parameter the level cannot pin,
$\sigma_\varepsilon$. The steady state is identical across all baseline
experiments.
